% Chapter 4. Source pages PDF 91--114 (printed 83--106).
\section{Definitions and examples}
\begin{definition}[Representable covariant functor]
\label{def:lbc-representable-cov}
\source{leinster-basic-category-theory:page-0092}
\dcref{4.1.1}
Let $\mathcal A$ be locally small. For $A\in\mathcal A$, the functor
$H^A=\mathcal A(A,-):\mathcal A\to\Set$ sends $B$ to $\mathcal A(A,B)$ and
$f:B\to B'$ to postcomposition $f(-)$.
\uses{def:lbc-locally-small,def:lbc-functor}
\end{definition}

\begin{remark}[Representable notation]
\label{rem:lbc-representable-notation}
\source{leinster-basic-category-theory:page-0092}
\dcref{4.1.2}
Local smallness ensures that $H^A(B)=\mathcal A(A,B)$ is a set. The action on
$g:B\to B'$ may be written $g\circ-$ or $g_*$.
\uses{def:lbc-representable-cov,def:lbc-locally-small}
\end{remark}

\begin{definition}[Representable functor]
\label{def:lbc-representable}
\source{leinster-basic-category-theory:page-0092}
\dcref{4.1.3}
A functor $X:\mathcal A\to\Set$ is representable if there are $A\in\mathcal A$
and a natural isomorphism $X\cong H^A$.
\uses{def:lbc-representable-cov,def:lbc-natural-iso}
\end{definition}

\begin{lemma}[Representability from a left adjoint]
\label{lem:lbc-representable-left-adjoint}
\source{leinster-basic-category-theory:page-0094}
\dcref{4.1.10}
Let $F:\mathcal A\to\mathcal B$ be left adjoint to $G:\mathcal B\to\mathcal A$
and let $A\in\mathcal A$. Then the functor
$\mathcal A(A,G(-)):\mathcal B\to\Set$ is represented by $F(A)$.
\uses{def:lbc-adjunction,def:lbc-representable-cov}
\end{lemma}
\begin{proof}
The adjunction bijection $\mathcal A(A,G(B))\cong\mathcal B(F(A),B)$ is
natural in $B$, which is precisely a natural isomorphism of the two set-valued
functors.
\end{proof}

\begin{proposition}[Left adjoint implies representable]
\label{prop:lbc-left-adjoint-representable}
\source{leinster-basic-category-theory:page-0095}
\dcref{4.1.11}
Any set-valued functor with a left adjoint is representable.
\uses{lem:lbc-representable-left-adjoint}
\end{proposition}
\begin{proof}
Apply Lemma~\ref{lem:lbc-representable-left-adjoint} to the unique object of
$\Set$; since $\mathcal A(1,-)$ is naturally isomorphic to the underlying-set
functor, the given functor is represented by the image of $1$ under its left
adjoint.
\end{proof}

\begin{definition}[Contravariant assembly of representables]
\label{def:lbc-representables-assembly}
\source{leinster-basic-category-theory:page-0096}
\dcref{4.1.15}
The functor $H^{\bullet}:\mathcal A^{\op}\to[\mathcal A,\Set]$ sends $A$ to
$H^A=\mathcal A(A,-)$ and $f$ to precomposition by $f$.
\uses{def:lbc-representable-cov,con:lbc-functor-category}
\end{definition}

\begin{definition}[Covariant representable in the opposite direction]
\label{def:lbc-representable-op}
\source{leinster-basic-category-theory:page-0096}
\dcref{4.1.16}
For locally small $\mathcal A$ and $A\in\mathcal A$, the functor
$H_A=\mathcal A(-,A):\mathcal A^{\op}\to\Set$ sends $g:B'\to B$ to
precomposition $p\mapsto p g$.
\uses{def:lbc-op-category,def:lbc-locally-small}
\end{definition}

\begin{definition}[Representable contravariant functor]
\label{def:lbc-representable-contra}
\source{leinster-basic-category-theory:page-0097}
\dcref{4.1.17}
For $A\in\mathcal A$, put $H_A=\mathcal A(-,A):\mathcal A^{\op}\to\Set$.
A presheaf is representable when it is naturally isomorphic to some $H_A$.
\uses{def:lbc-presheaf,def:lbc-representable-op}
\end{definition}

\begin{definition}[Yoneda embedding]
\label{def:lbc-yoneda-embedding}
\source{leinster-basic-category-theory:page-0098}
\dcref{4.1.21}
The Yoneda embedding is the functor
$H_{\bullet}:\mathcal A\to[\mathcal A^{\op},\Set]$ with $A\mapsto H_A$ and
$f:A\to A'$ acting by postcomposition.
\uses{def:lbc-representable-contra,con:lbc-functor-category}
\end{definition}

\begin{definition}[Generalized element]
\label{def:lbc-generalized-element}
\source{leinster-basic-category-theory:page-0100}
\dcref{4.1.25}
For a functor $X:\mathcal A^{\op}\to\Set$, a generalized element of $X$ of
shape $A$ is an element of $X(A)$; a universal element is a pair $(A,x)$ such
that every $(B,y)$ has a unique map $f:B\to A$ with $X(f)(x)=y$.
\uses{def:lbc-presheaf}
\end{definition}

\begin{definition}[Hom bifunctor]
\label{def:lbc-hom-bifunctor}
\source{leinster-basic-category-theory:page-0098}
\dcref{4.1.22}
For locally small $\mathcal A$, the functor
$\Hom_{\mathcal A}:\mathcal A^{\op}\times\mathcal A\to\Set$ sends $(A,B)$ to
$\mathcal A(A,B)$ and $(f:A'\to A,g:B\to B')$ to
$p\mapsto gpf$.
\uses{def:lbc-product-category,def:lbc-op-category,def:lbc-functor}
\end{definition}

\begin{remark}[The Hom bifunctor]
\label{rem:lbc-hom-bifunctor}
\source{leinster-basic-category-theory:page-0099}
\dcref{4.1.23}
The bifunctor $\Hom_{\mathcal A}$ packages the covariant representables into
$H^{\bullet}:\mathcal A^{\op}\to[\mathcal A,\Set]$; equivalently, it is the
currying of this functor under the product/functor-category adjunction.
\uses{def:lbc-hom-bifunctor,def:lbc-representables-assembly}
\end{remark}

\begin{remark}[Naturality of adjunctions]
\label{rem:lbc-representable-naturality}
\source{leinster-basic-category-theory:page-0099}
\dcref{4.1.24}
The naturality required in an adjunction is exactly the compatibility encoded by
the Hom bifunctor and its two representable-variable descriptions.
\uses{def:lbc-hom-bifunctor,def:lbc-adjunction}
\end{remark}

\section{The Yoneda lemma}
\begin{theorem}[Yoneda lemma]
\label{thm:lbc-yoneda}
\source{leinster-basic-category-theory:page-0102}
\dcref{4.2.1}
For locally small $\mathcal A$, $A\in\mathcal A$, and
$X:\mathcal A^{\op}\to\Set$, there is a natural bijection
\[\operatorname{Nat}(H_A,X)\cong X(A),\qquad \alpha\longmapsto\alpha_A(1_A).\]
The inverse sends $x\in X(A)$ to $\alpha^x_B(f)=X(f)(x)$.
\uses{def:lbc-yoneda-embedding,def:lbc-generalized-element}
\end{theorem}
\begin{proof}
For natural $\alpha$, naturality at $f:B\to A$ gives
$\alpha_B(f)=X(f)(\alpha_A(1_A))$, so $\alpha$ is determined by its displayed
element. Conversely, functoriality gives
$X(g)(X(f)(x))=X(fg)(x)$, proving naturality of $\alpha^x$; the two
constructions are inverse by evaluating at identities.
\end{proof}

\section{Consequences}
\begin{notation}[Invertible arrows]
\label{not:lbc-invertible-arrow}
\source{leinster-basic-category-theory:page-0107}
\dcref{4.3.1}
An arrow decorated by a tilde denotes an isomorphism.
\uses{def:lbc-isomorphism}
\end{notation}
\begin{corollary}[A representation is a universal element]
\label{cor:lbc-universal-element}
\source{leinster-basic-category-theory:page-0107}
\dcref{4.3.2}
Let $X:\mathcal A^{\op}\to\Set$. A representation of $X$ consists of an
object $A$ and an element $u\in X(A)$ such that, for every $B$ and
$x\in X(B)$, there is a unique $\bar x:B\to A$ with $X(\bar x)(u)=x$.
\uses{thm:lbc-yoneda,def:lbc-generalized-element}
\end{corollary}
\begin{proof}
By Yoneda, $u$ corresponds to a natural transformation $\widetilde u:H_A\to X$.
Its component at $B$ is bijective exactly when every $x\in X(B)$ has a unique
preimage $\bar x$, which is the displayed condition.
\end{proof}

\begin{corollary}[A covariant representation is a universal element]
\label{cor:lbc-universal-element-cov}
\source{leinster-basic-category-theory:page-0108}
\dcref{4.3.3}
Let $X:\mathcal A\to\Set$. A representation of $X$ consists of an object $A$
and an element $u\in X(A)$ such that, for every $B$ and $x\in X(B)$, there is
a unique $\bar x:A\to B$ with $X(\bar x)(u)=x$.
\uses{thm:lbc-yoneda,def:lbc-generalized-element}
\end{corollary}
\begin{proof}
This is the dual of Corollary~\ref{cor:lbc-universal-element}.
\end{proof}

\begin{corollary}[Yoneda embedding]
\label{cor:lbc-yoneda-full-faithful}
\source{leinster-basic-category-theory:page-0110}
\dcref{4.3.7}
For any locally small category $\mathcal A$, the Yoneda embedding
$H_{\bullet}:\mathcal A\to[\mathcal A^{\op},\Set]$ is full and faithful.
\uses{thm:lbc-yoneda,def:lbc-yoneda-embedding}
\end{corollary}
\begin{proof}
For $A,A'$, Yoneda identifies natural transformations $H_A\to H_{A'}$ with
elements of $H_{A'}(A)=\mathcal A(A,A')$, and the resulting map is precisely
postcomposition by the corresponding arrow.
\end{proof}

\begin{lemma}[Isomorphisms under full faithful functors]
\label{lem:lbc-full-faithful-nat}
\source{leinster-basic-category-theory:page-0111}
\dcref{4.3.8}
Let $J:\mathcal A\to\mathcal B$ be full and faithful and let $A,A'\in\mathcal A$.
(a) A map $f$ in $\mathcal A$ is an isomorphism iff $J(f)$ is an isomorphism.
(b) Every isomorphism $g:J(A)\to J(A')$ is $J(f)$ for a unique isomorphism
$f:A\to A'$. (c) $A\cong A'$ iff $J(A)\cong J(A')$.
\uses{def:lbc-full-faithful,def:lbc-isomorphism}
\notready
\end{lemma}
\begin{proof}
The source defers this proof to Exercise 4.3.15.
\end{proof}

\begin{corollary}[Isomorphic representables]
\label{cor:lbc-representable-iso}
\source{leinster-basic-category-theory:page-0112}
\dcref{4.3.10}
For locally small $\mathcal A$ and $A,A'\in\mathcal A$,
$H_A\cong H_{A'}\iff A\cong A'\iff H^A\cong H^{A'}$.
\uses{cor:lbc-yoneda-full-faithful,lem:lbc-full-faithful-nat}
\end{corollary}
\begin{proof}
The first equivalence follows from the full faithfulness in Corollary
\ref{cor:lbc-yoneda-full-faithful} and part (c) of Lemma~\ref{lem:lbc-full-faithful-nat};
the second is its dual.
\end{proof}
